\label{sec:results}

\subsection{\textit{Allen v.\ Milligan} (2023): \textit{Gingles} Reaffirmed}

\textit{Allen v.\ Milligan}\footnote{\textit{Allen v.\ Milligan}, 599 U.S.\
1 (2023).} is the most important post-\textit{Shelby} Supreme Court decision
on redistricting and VRA Section~2. The case arose from Alabama's 2021
congressional redistricting. Alabama has seven congressional districts;
the NAACP Legal Defense Fund and other plaintiffs argued that Alabama's
plan violated Section~2 by packing Black voters into a single majority-Black
district (the 7th) while cracking the remaining Black population across
six predominantly white districts, even though Black Alabamans constitute
approximately 27\% of the state's population.

The district court (a three-judge court under 52~U.S.C.~\S~10304(a))
held that plaintiffs had demonstrated a likelihood of success on the merits
and ordered Alabama to draw a new map with a second district in which
Black voters had an opportunity to elect their preferred candidates.
Alabama appealed; the Supreme Court issued a stay pending appeal.

The Supreme Court affirmed by a 5--4 vote.\footnote{Chief Justice Roberts
joined the four-Justice liberal bloc in the majority; Justice Kavanaugh
wrote a concurrence.} Chief Justice Roberts, writing for the majority,
reaffirmed the \textit{Gingles} framework as the operative standard for
Section~2 redistricting claims. He rejected Alabama's arguments that
Section~2 required the drawing of majority-minority districts only to be
evaluated through a race-neutral analysis, or that the Equal Protection
Clause's colorblindness requirement precluded a court from considering
race in evaluating whether a second majority-minority district was required.

The majority's core holding: \textit{Gingles} means what it says. Where
the three preconditions are met (geographic concentration and compactness,
political cohesion, majority bloc voting), the totality of circumstances
analysis must be conducted, and if it shows vote dilution, a remedial
majority-minority district must be drawn. The Court reaffirmed that using
race as a factor in redistricting to comply with the VRA does not
automatically violate the Equal Protection Clause---it is the use of race
to \textit{exceed} VRA requirements that triggers strict scrutiny.

\textit{Allen v.\ Milligan} matters for the post-\textit{Shelby} landscape
in two ways. First, it confirms that Section~2 vote-dilution claims for
redistricting remain viable and that courts will order remedial maps
when plaintiffs prevail. Second, and more subtly, it demonstrates the
power of the \textit{Gingles} framework's compactness analysis: the
district court found, and the Supreme Court affirmed, that Alabama's
Black population was sufficiently concentrated to constitute a majority
in a second congressional district---a finding driven by geographic and
demographic analysis of the kind that algorithmic redistricting tools
like \textsc{bisect} can perform systematically.

\subsection{State VRA Analogs}

Several states have enacted state-law analogs to the VRA that provide
additional protection beyond what federal law requires post-\textit{Shelby}.

\textbf{California Voting Rights Act (CVRA)}: Enacted in 2001 (Cal.\ Elec.\
Code \S\S\ 14025--14032), the CVRA prohibits at-large elections that dilute
minority voting rights and applies a lower standard than the federal
\textit{Gingles} test: it does not require a showing of geographically
compact majority-minority district feasibility. The CVRA applies to local
governments and school districts; it has generated substantial litigation
and has resulted in numerous jurisdictions converting from at-large to
by-district election systems.

\textbf{New York Voting Rights Act (NYVRA)}: Enacted in 2022 as part of
the John R.\ Lewis Voting Rights Act of New York (S51001-A), the NYVRA
provides the most comprehensive state-level VRA analog in the country.
It covers all levels of government, employs a results standard similar
to the federal Section~2, and includes enhanced remedies. Critically,
the NYVRA restores a preclearance-like mechanism for jurisdictions with
recent voting rights violations, requiring them to obtain approval from
the New York Attorney General before implementing changes to voting laws.
This represents the most direct state-level substitute for federal Section~5
preclearance enacted anywhere in the country.

\textbf{Washington Voting Rights Act (WVRA)}: Enacted in 2018 (Wash.\ Rev.\
Code \S\ 29A.92), the WVRA prohibits dilution of minority voting rights in
local elections and provides a results-based standard that applies broadly
to at-large districts.

These state analogs vary in their applicability to congressional redistricting
(which is governed primarily by federal law) vs.\ state and local redistricting
(where states have broader authority). The NYVRA is the only state VRA analog
that includes a preclearance-like mechanism.

\subsection{The Enforcement Asymmetry in Practice}

Table~\ref{tab:enforcement} summarizes the enforcement asymmetry between
the pre-\textit{Shelby} and post-\textit{Shelby} environments.

\begin{table}[ht]
\centering
\caption{VRA Redistricting Enforcement: Pre-\textit{Shelby} vs.\ Post-\textit{Shelby}}
\label{tab:enforcement}
\begin{tabular}{lll}
\toprule
Dimension & Pre-\textit{Shelby} (through 2013) & Post-\textit{Shelby} (2013--present) \\
\midrule
Primary tool & Section~5 preclearance & Section~2 litigation \\
Timing & Ex ante (before enactment) & Ex post (after enactment) \\
Process & Administrative (DOJ) or D.C.\ court & Federal district court \\
Burden & State proves non-retrogression & Plaintiff proves dilution \\
Timeline & Weeks to months & Years \\
Covered states & 9 fully + parts of others & None (Section~4(b) invalid) \\
Standard & Non-retrogression & \textit{Gingles} + totality \\
\bottomrule
\end{tabular}
\end{table}

The enforcement asymmetry creates a strategic environment in which former
covered states can enact potentially discriminatory plans knowing that
any legal challenge will take years to resolve. The 2012 Texas congressional
plan, challenged immediately after enactment, was not replaced until after
the 2014 and 2016 elections; Alabama's 2021 plan was not remedied until
after the 2022 election (following \textit{Allen v.\ Milligan}). In a
world with functioning Section~5 preclearance, neither plan would have
been implemented without federal approval.
